\begin{thaiabstract}
การทำแบบสอบถามที่ผู้ป่วยกรอกเองมีข้อจำกัดด้านการเข้าถึง โดยเฉพาะสำหรับผู้สูงอายุและผู้ป่วยที่มีข้อจำกัดทางร่างกาย 
จึงจำเป็นต้องให้พยาบาลเป็นผู้สัมภาษณ์ ซึ่งต้องใช้บุคลากรและเวลามาก
วิทยานิพนธ์นี้นำเสนอตัวแทนสนทนาด้วยเสียงภาษาไทยสำหรับเก็บข้อมูลผลลัพธ์ทางสุขภาพ
ที่รายงานโดยผู้ป่วยในรูปแบบอิเล็กทรอนิกส์
เพื่อลดภาระของผู้ป่วยและบุคลากรทางการแพทย์
ระบบรองรับแบบประเมิน 3 แบบ ได้แก่ DASH และ TESS สำหรับรยางค์บนและรยางค์ล่าง
โดยจะถามคำถามแต่ละข้อด้วยเสียงที่สร้างไว้ล่วงหน้าพร้อมภาพประกอบ
แปลงคำตอบของผู้ป่วยเป็นข้อความด้วยระบบรู้จำเสียงพูดอัตโนมัติ
และใช้แบบจำลองภาษาขนาดใหญ่ในการประมวลผลและจับคู่คำตอบกับตัวเลือกของแต่ละข้อ
โดยทุกองค์ประกอบทำงานบนเซิร์ฟเวอร์ภายในเพื่อรักษาความเป็นส่วนตัวของข้อมูลผู้ป่วย
และได้มีการประเมินความพร้อมสำหรับการใช้งานทางคลินิก 
ด้วยการวัดความสอดคล้องระหว่างคะแนนจากระบบกับคะแนนจากการสัมภาษณ์โดยพยาบาล ทั้งในระดับรายข้อและระดับคะแนนรวม

% TODO: เพิ่มผลการทดลองโดยสรุป เช่น ค่าความสอดคล้องรายข้อ (weighted kappa) และคะแนนรวม (ICC)
\end{thaiabstract}
